\section{实射影空间}

在本节中我们需要不断地引用有关商空间的概念和结果（第 I 章，§ 3），特别是下面两条性质，为方便起见我们把它们陈述为引理：

设 E 是拓扑空间，R 是 E 上的一个等价关系，A 是 E 的一个子集，\(R_{A}\) 是 R 在 A 上诱导的等价关系，f 是 E 到 E/R 上的典范映射。那么：

引理 1. 若 A 中每个关于 \(R_{A}\) 饱和的开（相应地，闭）集，都是 E 中关于 R 饱和的开（相应地，闭）集在 A 上的迹，则商空间 \(A/R_{A}\) 同胚于 E/R 的子空间 \(f(A)\)。特别地，当 A 在 E 中是开的或闭的且关于 R 饱和时，情形就是如此。

这由第 I 章，§ 3，第 6 小节，命题 10 得出。

引理 2. 若存在连续映射 \(g\) ——\(\mathbf{E}\) 到 \(\mathbf{A}\) 上的——，使得对一切 \(x \in \mathbf{E}\) ，\(g(x)\) 属于 \(x\) 的等价类，则商空间 \(\mathrm{A} / \mathrm{R}_{\mathbf{A}}\) 同胚于 \(\mathbf{E} / \mathbf{R}\)。

这就是第 I 章，§ 3，第 6 小节，命题 10 的推论 2。

